\chapter{The Inclusion--Exclusion Principle}

\section{Overview}

The goal of this project is to formalize the inclusion--exclusion principle
for finite sets in Lean 4 using Mathlib.

Given a finite indexed family of finite sets
\[
\{S_b\}_{b \in s},
\]
the inclusion--exclusion principle expresses the cardinality of the union
\[
\bigcup_{b \in s} S_b
\]
as an alternating sum over intersections of subfamilies.

We formalize this statement using finite sets (`Finset`) and prove the result
by induction on the indexing set.

The full formal statement is given in Lean as:

\lean{InclusionExclusion.lean, inclusion_exclusion}

---

\section{Basic Union Identities}

We begin with elementary cardinality results for unions and intersections.

\begin{lemma}[Two-set union formula]\label{card_union_two}
For finite sets \(A,B\),
\[
|A \cup B|
=
|A| + |B| - |A \cap B|.
\]
\end{lemma}

\lean{Union.lean, card_union_two}

---

\begin{lemma}[Disjoint union]\label{disjoint_union}
If \(A\) and \(B\) are disjoint, then
\[
|A \cup B|
=
|A| + |B|.
\]
\end{lemma}

\lean{Union.lean, disjoint_union}

---

\begin{lemma}[Idempotent union]\label{same_union}
For every finite set \(A\),
\[
|A \cup A| = |A|.
\]
\end{lemma}

\lean{Union.lean, same_union}

---

\section{Finite Indexed Unions}

We next establish structural lemmas for finite indexed unions.

\begin{lemma}\label{biUnion_insert}
Let \(s\) be a finite index set and \(b \notin s\). Then
\[
\bigcup_{x \in \{b\} \cup s} S_x
=
S_b \cup \bigcup_{x \in s} S_x.
\]
\end{lemma}

\lean{InclusionExclusion.lean, biUnion_insert}

---

\begin{lemma}\label{biUnion_empty}
\[
\bigcup_{x \in \varnothing} S_x = \varnothing.
\]
\end{lemma}

\lean{InclusionExclusion.lean, biUnion_empty}

---

\begin{lemma}\label{mem_biUnion}
For \(x\),
\[
x \in \bigcup_{b \in s} S_b
\iff
\exists b \in s,\ x \in S_b.
\]
\end{lemma}

\lean{InclusionExclusion.lean, mem_biUnion}

---

\section{Finite Intersections}

We interpret intersections as infima in the lattice structure on `Finset`.

\begin{lemma}\label{mem_inf}
For finite \(T\),
\[
x \in \bigcap_{b \in T} S_b
\iff
\forall b \in T,\ x \in S_b.
\]
\end{lemma}

\lean{InclusionExclusion.lean, mem_inf_iff_forall}

---

\section{Powersets and Alternating Sums}

We now study subsets of finite sets and alternating combinatorial sums.

\begin{lemma}\label{mem_powerset}
For finite sets \(T,s\),
\[
T \in \mathcal P(s)
\iff
T \subseteq s.
\]
\end{lemma}

\lean{InclusionExclusion.lean, mem_powerset}

---

\begin{lemma}\label{alternating_sum_powerset}
If \(s \neq \varnothing\), then
\[
\sum_{T \subseteq s} (-1)^{|T|}
=
0.
\]
\end{lemma}

\lean{InclusionExclusion.lean, alternating_sum_powerset}

---

\section{The Inclusion--Exclusion Principle}

We now prove the main theorem.

\begin{theorem}[Inclusion--Exclusion]\label{inclusion_exclusion}
Let \(s\) be a finite index set and let
\[
S : \beta \to \mathrm{Finset}(\alpha)
\]
be a finite family of finite sets. Then
\[
\left|
\bigcup_{b \in s} S_b
\right|
=
\sum_{\varnothing \neq T \subseteq s}
(-1)^{|T|+1}
\left|
\bigcap_{b \in T} S_b
\right|.
\]
\end{theorem}

\lean{InclusionExclusion.lean, inclusion_exclusion}

---

\section{Proof Strategy}

The proof proceeds by induction on the finite index set \(s\),
using:
\begin{itemize}
\item the two-set union formula (\ref{card_union_two}),
\item structural lemmas on `biUnion`,
\item and combinatorial cancellation in alternating sums over subsets.
\end{itemize}